\documentclass[12pt,a4paper]{article}

% -------------------------------------------------
% Sprache / Kodierung
% -------------------------------------------------
\usepackage[T1]{fontenc}
\usepackage[utf8]{inputenc}
\usepackage[ngerman]{babel}
\usepackage{lmodern}
\usepackage{microtype}

% -------------------------------------------------
% Mathematik
% -------------------------------------------------
\usepackage{amsmath,amssymb,amsthm,mathtools}

% -------------------------------------------------
% Layout / Grafik / Farben
% -------------------------------------------------
\usepackage[a4paper,margin=2.3cm]{geometry}
\usepackage{booktabs}
\usepackage{xcolor}
\usepackage[hidelinks,unicode]{hyperref}
\pdfstringdefDisableCommands{%
  \def\ü{ue}\def\ö{oe}\def\ä{ae}\def\Ü{Ue}\def\Ö{Oe}\def\Ä{Ae}\def\ss{ss}%
}
\usepackage[most]{tcolorbox}
\usepackage{enumitem}
\usepackage{array}

% -------------------------------------------------
% Farben (konsistent mit Vorgängerdokumenten)
% -------------------------------------------------
\definecolor{lückenrot}{RGB}{180,40,40}
\definecolor{bewiesengrün}{RGB}{30,100,50}
\definecolor{warnorange}{RGB}{200,100,10}
\definecolor{hintergrundblau}{RGB}{235,242,250}
\definecolor{hintergrundgelb}{RGB}{255,252,225}
\definecolor{adischviolett}{RGB}{90,20,140}
\definecolor{fehlerrot}{RGB}{200,0,0}
\definecolor{konjviolett}{RGB}{80,0,130}
\definecolor{e8grau}{RGB}{55,55,90}
\definecolor{ikobraun}{RGB}{120,70,20}

% -------------------------------------------------
% tcolorbox-Stile
% -------------------------------------------------
\tcbset{
  lücke/.style={colback=red!8, colframe=lückenrot, fonttitle=\bfseries,
    title={Offene Lücke}},
  ergebnis/.style={colback=green!7, colframe=bewiesengrün,
    fonttitle=\bfseries, title={Bewiesenes Ergebnis}},
  warnung/.style={colback=orange!10, colframe=warnorange,
    fonttitle=\bfseries, title={Achtung / Heuristik}},
  adisch/.style={colback=violet!6, colframe=adischviolett,
    fonttitle=\bfseries, title={2-adische Perspektive}},
  kernidee/.style={colback=hintergrundblau, colframe=blue!60!black,
    fonttitle=\bfseries, title={Kernidee}},
  konj/.style={colback=violet!5, colframe=konjviolett, fonttitle=\bfseries,
    title={Konjektur (unbewiesen)}},
  lte/.style={colback=yellow!8, colframe=e8grau, fonttitle=\bfseries,
    title={LTE-Analyse}},
  hauptsatz/.style={colback=green!10, colframe=bewiesengrün!80!black,
    fonttitle=\bfseries\large, title={Stärkster beweisbarer Satz}},
  analogie/.style={colback=brown!6, colframe=ikobraun,
    fonttitle=\bfseries, title={Geometrische Analogie}},
  ikosaeder/.style={colback=blue!4, colframe=blue!45!black,
    fonttitle=\bfseries, title={Ikosaeder-Perspektive}},
}

% -------------------------------------------------
% Theorem-Umgebungen
% -------------------------------------------------
\newtheorem{definition}{Definition}[section]
\newtheorem{proposition}[definition]{Proposition}
\newtheorem{theorem}[definition]{Theorem}
\newtheorem{lemma}[definition]{Lemma}
\newtheorem{korollar}[definition]{Korollar}
\newtheorem{bemerkung}[definition]{Bemerkung}
\newtheorem{hypothese}[definition]{Hypothese}
\newtheorem{satz}[definition]{Satz}

\newenvironment{beweis}{\begin{proof}[Beweis]}{\end{proof}}

% -------------------------------------------------
% Makros
% -------------------------------------------------
\newcommand{\vii}{\nu_2}
\newcommand{\U}{\mathcal{U}}
\newcommand{\N}{\mathbb{N}}
\newcommand{\Z}{\mathbb{Z}}
\newcommand{\R}{\mathbb{R}}
\newcommand{\E}{\mathbb{E}}
\newcommand{\Ztwo}{\mathbb{Z}_2}
\newcommand{\abs}[1]{\lvert#1\rvert}
\newcommand{\atwo}[1]{\lvert#1\rvert_2}
\newcommand{\logF}{\log F(B)}
\newcommand{\ZmZ}{\Z/12\Z}
\newcommand{\Ssph}{\mathbb{S}^2}

% -------------------------------------------------
% Dokument
% -------------------------------------------------
\title{%
  Schütte-Ikosaeder-Spannung und Collatz-Dämpfung:\\[0.4em]
  Geometrische Analogie und ihre Grenzen\\[0.5em]
  \large Ergänzung zu \textit{collatz\_lyapunov\_konditioniert.tex}%
}
\author{Thomas Hoffbauer}
\date{16.\ Juni 2026}

\begin{document}
\maketitle

\begin{abstract}
Dieses Dokument untersucht, ob die \emph{Schütte-Ikosaeder-Spannung} ---
im Projektkontext das topologische Maß der Instabilität in Schüttes Beweis
des 12-Kuss-Theorems von 1953 --- zur Verstärkung des negativen
Lyapunov-Exponenten der Collatz-Dämpfung beitragen kann.

Der mathematische Kern der Analogie ist folgende Beobachtung: Das
Ikosaeder hat genau \textbf{12 Ecken}; die Gruppe $\ZmZ$ hat genau
\textbf{12 Restklassen}; davon sind genau \textbf{4} die EABC-Klassen,
die die Collatz-Dynamik auf ungeraden Zahlen tragen.

Die Analyse trennt konsequent drei Ebenen:
\begin{enumerate}[noitemsep]
  \item \textcolor{bewiesengrün}{\textbf{Bewiesen}}: Das 12-Kuss-Theorem
        (Schütte/van der Waerden 1953); die EABC-Klassenstruktur; die
        LTE-Schranke.
  \item \textcolor{warnorange}{\textbf{Geometrische Analogie}}: Die
        strukturelle Parallele zwischen Ikosaeder-Ecken und Restklassen
        mod~12; eine Übertragung der Spannungsungleichung auf die
        EABC-Übergangsmatrix.
  \item \textcolor{lückenrot}{\textbf{Offen}}: Ergodizität der Collatz-Bahn;
        Zyklenfreiheit; ein direkter quantitativer Beitrag der
        Ikosaeder-Geometrie zum Lyapunov-Exponenten.
\end{enumerate}

\medskip\noindent\textbf{Fazit:}
Die Ikosaeder-Spannung liefert eine kohärente geometrische Sprache für die
Endlichkeit und Vollständigkeit des EABC-Zustandsraums, aber keinen
direkten Beweis der Dämpfung.
\end{abstract}

\tableofcontents
\bigskip

\begin{tcolorbox}[warnung, title={Einordnung dieses Dokuments}]
Die Collatz-Vermutung ist unbewiesen.
Alle Drift-Rechnungen aus \textit{collatz\_lyapunov\_konditioniert.tex}
sind heuristisch.
Dieses Dokument fügt eine geometrische Perspektive hinzu, die in sich
mathematisch korrekt ist, aber den Status des Hauptarguments nicht verändert.
\end{tcolorbox}

% ====================================================
\section{Schütte-Ikosaeder-Spannung: Mathematischer Gehalt}
\label{sec:schuette}
% ====================================================

\subsection{Das 12-Kuss-Theorem}

\begin{tcolorbox}[ergebnis, title={Satz von Schütte und van der Waerden (1953) -- bewiesen}]
\textbf{Theorem (12-Kuss-Theorem):}
In $\R^3$ können maximal $12$ nicht-überlappende Einheitskugeln eine
zentrale Einheitskugel gleichzeitig berühren.
Äquivalent: Es gibt maximal $12$ Einheitsvektoren $x_1,\dots,x_n \in \Ssph$
mit paarweisem Winkel
\[
  \angle(x_i, x_j) \;\geq\; 60°
  \quad \Longleftrightarrow \quad
  \langle x_i, x_j \rangle \;\leq\; \tfrac{1}{2}
  \qquad \text{für alle } i \neq j.
\]
\end{tcolorbox}

Der Beweis von Schütte und van der Waerden (1953) ist bemerkenswert, weil
er \emph{knapp} ist: Tatsächlich können die 12 Kugeln so angeordnet werden,
dass jede leicht beweglich bleibt --- es gibt kein eindeutiges starres
Optimum wie in Dimension~8 (Ikosaeder ist nicht die einzige optimale
Konfiguration).
Genau diese Restfreiheit nennt das Projekt \textbf{Schütte-Spannung}.

\begin{definition}[Schütte-Spannung, nach \textit{Final RH.tex}]
\label{def:spannung}
Die \emph{Schütte-Spannung} $\sigma \geq 0$ ist das Maß der
topologischen Instabilität im $\R^3$-Kissing-Problem:
\[
  \sigma \;:=\; 1 - \frac{\text{Kissing-Zahl}(\R^3)}{%
  \text{Kissing-Zahl}(\R^8)} \;=\; 1 - \frac{12}{240} \;=\; \frac{19}{20}.
\]
In $\R^8$ verschwindet $\sigma$ vollständig: Das $E_8$-Gitter realisiert
die Kissing-Zahl $240$ in einer einzigartigen, starren Konfiguration
(Viazovska, 2022 Fields-Medaille).
\end{definition}

\begin{bemerkung}[Gregory-Newton-Lücke]
Die Schütte-Spannung manifestiert sich geometrisch in der
\emph{Gregory-Newton-Lücke}: Zwischen der 12-Kugel-Konfiguration und dem
Aufbrechen (ab 13 Kugeln) gibt es einen positiven Winkelspielraum von
etwa $\theta_\mathrm{GN} \approx 0.099$ Bogenmass~$\approx 5.7°$.
Die 12~Kugeln sitzen also \emph{nicht} in starrer Gleichgewichtslage,
sondern in einem \emph{gespannten}, aber beweglichen Zustand.
\end{bemerkung}

\subsection{Schütte-Spannung im Projektkontext}

Im vorliegenden Projekt taucht die Schütte-Spannung in mehreren Kontexten auf:

\begin{enumerate}[label=(\alph*)]
  \item \textbf{\textit{Final RH.tex}}: Schütte-Spannung $\sigma$ als Maß
        für den Unterschied zwischen der instabilen $\R^3$-Konfiguration
        (Kissing-Zahl $12$) und der rigiden $E_8$-Situation (Kissing-Zahl
        $240$). In $\R^3$ ist $\sigma > 0$; in $\R^8$ gilt $\sigma = 0$.

  \item \textbf{\textit{Pentagon Holonomie.tex}}: Die
        \emph{Schütte-Erweiterung} ist eine Folge diskreter
        Verfeinerungsschritte $\Sigma_n: \mathcal{T} \mapsto \mathcal{I}_n$
        vom Tetraederkern zur Mikrosaeder-Schale. Der Begriff ist strukturell
        gemeint: Die Spannung entsteht an jedem Ikosaeder/Dodekaeder-Umschlag.

  \item \textbf{\textit{Scholze\_12.tex}}: Die
        \emph{Schüttespannung $57$--$59$--$61$} als lokaler arithmetischer
        Resonanzmarker, der Kompositheit und Primheit in unmittelbarer
        Nachbarschaft zeigt.
\end{enumerate}

\begin{tcolorbox}[ikosaeder]
\textbf{Gemeinsamer Kern:}
Schütte-Spannung ist stets ein Maß für \emph{strukturelle Vollständigkeit
unter Zwangsbedingung}: Die Konfiguration ist maximal besetzt (12 Kugeln,
12 Restklassen, 12 Ecken), aber nicht starr --- der verbleibende
``Spielraum'' ist die Spannung.
\end{tcolorbox}

% ====================================================
\section{Die 12 Ecken des Ikosaeders und die 12 Restklassen mod 12}
\label{sec:zwolf}
% ====================================================

\subsection{Kombinatorik des Ikosaeders}

\begin{tcolorbox}[ergebnis, title={Ikosaeder-Grunddaten -- bewiesen (Euler, Euklid)}]
Das reguläre Ikosaeder $\mathcal{I}$ hat:
\begin{itemize}[noitemsep]
  \item $|V(\mathcal{I})| = 12$ Ecken (Vertices), je mit Grad $5$,
  \item $|E(\mathcal{I})| = 30$ Kanten,
  \item $|F(\mathcal{I})| = 20$ Flächen (gleichseitige Dreiecke).
\end{itemize}
Euler-Charakteristik: $12 - 30 + 20 = 2 = \chi(S^2)$.
Die Gram-Identität der Inzidenzmatrix (vgl.\ \textit{IkoDo\_transition.tex}):
\[
  B_I^\top B_I \;=\; 5\,I_{12} + A_I,
\]
wobei $A_I \in \{0,1\}^{12 \times 12}$ die Adjazenzmatrix ist.
\end{tcolorbox}

\subsection{Die 12 Restklassen mod 12 und ihre Rollen}

\begin{tcolorbox}[ergebnis, title={Struktursatz der Restklassen mod 12 -- bewiesen}]
Die Gruppe $\ZmZ = \{0, 1, 2, \dots, 11\}$ zerfällt in drei disjunkte
Typen:
\begin{center}
\renewcommand{\arraystretch}{1.3}
\begin{tabular}{lll}
\toprule
Typ & Restklassen & Rolle in Collatz \\
\midrule
EABC (aktiv) & $\{1, 5, 7, 11\}$ & ungerade, $\gcd(\cdot, 6) = 1$ \\
\quad E & $\{1\}$ & $n \equiv 1 \pmod{12}$ \\
\quad A & $\{5\}$ & $n \equiv 5 \pmod{12}$ \\
\quad B & $\{7\}$ & $n \equiv 7 \pmod{12}$ \\
\quad C & $\{11\}$ & $n \equiv 11 \pmod{12}$ \\
\midrule
2-Drain (passiv) & $\{0, 2, 4, 6, 8, 10\}$ & gerade, $2 \mid n$ \\
3-Drain (passiv) & $\{3, 9\}$ & $3 \mid n$, $2 \nmid n$ \\
\bottomrule
\end{tabular}
\end{center}
Von den 12 Restklassen sind also genau $\mathbf{4}$ EABC-aktiv und
$\mathbf{8}$ inaktiv.
\end{tcolorbox}

\subsection{Die strukturelle Parallele}

\begin{tcolorbox}[analogie]
\textbf{Beobachtung (Analogie, kein Beweis):}
Das Ikosaeder hat 12 Ecken; $\ZmZ$ hat 12 Elemente.
In beiden Fällen gibt es eine \emph{4-elementige aktive Teilmenge}:
\begin{itemize}[noitemsep]
  \item Ikosaeder: 4 ausgezeichnete Ecken einer eingebetteten Tetraeder-Untergruppe
        (der reguläre Tetraeder passt in das Ikosaeder ein, seine 4 Ecken
        bilden eine maximal separierte Teilmenge).
  \item Mod 12: Die 4 EABC-Klassen $\{1, 5, 7, 11\}$ sind genau die
        Einheiten der Gruppe $(\Z/12\Z)^\times$ --- maximal separiert im
        Sinne: kein Element ist durch 2 oder 3 teilbar.
\end{itemize}

\medskip
\textbf{Numerische Entsprechung:}
\[
  \begin{array}{c|c}
  \text{Ikosaeder} & \ZmZ \\ \hline
  12 \text{ Ecken} & 12 \text{ Restklassen} \\
  4 \text{ Tetraeder-Ecken} & 4 \text{ EABC-Klassen} \\
  8 \text{ ``äußere'' Ecken} & 8 \text{ inaktive Klassen}
  \end{array}
\]
\end{tcolorbox}

\begin{bemerkung}[Winkeldefekt]
\label{bem:winkeldefekt}
Im regulären Ikosaeder beträgt der Winkeldefekt an jeder Ecke:
\[
  \delta_{\text{Iko}} \;=\; 2\pi - 5 \cdot \frac{\pi}{3}
  \;=\; 2\pi - \frac{5\pi}{3}
  \;=\; \frac{\pi}{3} \;\approx\; 60°.
\]
Zum Vergleich: Das Tetraeder hat $\delta_T = 2\pi - 3 \cdot \frac{\pi}{3} = \pi$,
der Würfel hat $\delta_W = \frac{\pi}{2}$.
Der Winkeldefekt ist geometrisches Maß der ``Spannung'' an einer Ecke.
\end{bemerkung}

% ====================================================
\section{Übertragung der Spannungsungleichung auf die EABC-Übergangsmatrix}
\label{sec:transfer}
% ====================================================

\subsection{Die Kissing-Ungleichung}

Die Kernidee des Schütte-van der Waerden-Beweises ist eine Energieungleichung.
Für $n$ Einheitsvektoren $x_1, \dots, x_n \in \Ssph$ mit paarweisem
Innenprodukt $\langle x_i, x_j \rangle \leq c$ gilt:
\[
  0 \;\leq\; \left\|\sum_{i=1}^{n} x_i\right\|^2
  \;=\; n + 2 \sum_{i < j} \langle x_i, x_j \rangle
  \;\leq\; n + n(n-1)\,c.
\]
Diese triviale Ungleichung liefert allein noch keine obere Schranke für~$n$.
Schütte und van der Waerden kombinierten sie mit einer feineren
sphärisch-geometrischen Abschätzung (Winkelsummen auf $S^2$), die
tatsächlich $n \leq 12$ ergibt.

\begin{lemma}[Schütte-Ungleichung, vereinfacht]
\label{lem:schuette}
Sei $f: [-1, \tfrac{1}{2}] \to \R$ eine \emph{positiv definite}
Funktion auf $S^2$ (im Sinne von Schoenberg), etwa
\[
  f(t) \;=\; \bigl(1 - 2t\bigr)^{-1} \quad \text{(Modellwahl von Schütte)}.
\]
Dann gilt für $n$ Punkte $x_1, \dots, x_n \in S^2$ mit
$\langle x_i, x_j \rangle \leq \tfrac{1}{2}$ für $i \neq j$:
\[
  \sum_{i,j=1}^{n} f\!\bigl(\langle x_i, x_j \rangle\bigr) \;\geq\; 0.
\]
Ausgeschrieben ergibt sich (mit geeigneter $f$) eine obere Schranke
$n \leq 12$.
\end{lemma}

\begin{tcolorbox}[warnung, title={Wichtige Einschränkung}]
Die obige vereinfachte Formulierung des Schütte-Beweises ist qualitativ
korrekt, aber nicht vollständig. Der vollständige Beweis (Schütte und
van der Waerden, 1953) verwendet eine spezifische positiv-definite Funktion
und eine sorgfältige Normierungsrechnung.
\end{tcolorbox}

\subsection{Analogie: EABC-Übergangsmatrix}

Die odd-to-odd Collatz-Abbildung $\mathcal{U}(n) = (3n+1)/2^{\vii(3n+1)}$
induziert Übergänge zwischen den EABC-Klassen.
Auf der Ebene $\ZmZ$ ist die Dynamik allerdings \emph{nicht} rein durch
mod-12 bestimmt: $\mathcal{U}(n) \pmod{12}$ hängt von $\vii(3n+1)$ ab,
das wiederum von höheren Zweierpotenzen in~$n$ abhängt.

Wir betrachten deshalb die \emph{stochastische Approximation}: Definiere
die $(4 \times 4)$-Matrix $P_{\text{EABC}}$ durch
\[
  (P_{\text{EABC}})_{ab} \;:=\; P\bigl(\mathcal{U}(n) \equiv b
  \pmod{12} \;\big|\; n \equiv a \pmod{12}\bigr),
  \quad a, b \in \{1, 5, 7, 11\},
\]
wobei die Wahrscheinlichkeit über die stationäre Verteilung der Kettenlängen
$l$ gebildet wird.

\begin{tcolorbox}[ergebnis, title={Bekannte Fakten über $P_{\text{EABC}}$ -- bewiesen}]
\begin{enumerate}[noitemsep]
  \item $P_{\text{EABC}}$ ist eine \emph{substochastische} Matrix: Zeilensummen
        $\leq 1$, da ein Teil der Wahrscheinlichkeitsmasse in den
        2-Drain-Bereich abfließt.
  \item Der \emph{Spektralradius} $\rho(P_{\text{EABC}}) \leq 1$
        (Perron-Frobenius für nicht-negative Matrizen).
  \item Die Matrix ist endlich ($4 \times 4$): Der Zustandsraum ist
        abgeschlossen unter mod-12.
\end{enumerate}
\end{tcolorbox}

\begin{tcolorbox}[analogie]
\textbf{Spannungsanalogie (Analogie, kein Beweis):}

Das Kissing-Problem in $\R^3$ beschreibt, wie viele Vektoren $x_i \in S^2$
mit paarweisem Winkelabstand $\geq 60°$ existieren können --- maximal~12.
Übertragen auf $\ZmZ$:
\begin{itemize}[noitemsep]
  \item Die 12 Restklassen spielen die Rolle der 12 Kügelanordnungen.
  \item Die 4 EABC-Klassen sind die ``aktiven'' Kisser --- sie haben
        maximale gegenseitige modulare Distanz (kein EABC-Element ist
        Summe zweier anderer in $\Z/12\Z$ mit $\gcd = 1$).
  \item Die Kissing-Bedingung $n \leq 12$ entspricht dem Fakt, dass es
        \emph{maximal} 12 Restklassen mod~12 gibt --- keine 13.\ Klasse
        passt hinzu.
\end{itemize}

\medskip
\textbf{Strukturelle Parallele zur Spektralschranke:}
Falls die Schütte-Ungleichung (Lemma~\ref{lem:schuette}) auf eine
Matrixfunktion $f(P_{\text{EABC}})$ übertragbar wäre, ergäbe sich
\[
  \sum_{i,j} f\!\bigl((P_{\text{EABC}})_{ij}\bigr) \;\geq\; 0
  \quad \Longrightarrow \quad
  \rho(P_{\text{EABC}}) \;<\; 1.
\]
Dies ist jedoch \emph{nicht bewiesen}: Die Matrixeinträge von
$P_{\text{EABC}}$ sind keine Winkelkosinusse von Einheitsvektoren.
\end{tcolorbox}

\subsection{Was die Finiteness-Struktur tatsächlich liefert}

\begin{proposition}[Endlichkeit des Zustandsraums -- bewiesen]
\label{prop:endlich}
Der EABC-Zustandsraum $\{1, 5, 7, 11\} \subset \ZmZ$ ist endlich.
Insbesondere:
\begin{enumerate}[noitemsep]
  \item Die Übergangsmatrix $P_{\text{EABC}}$ ist eine endliche $(4\times 4)$-Matrix.
  \item Jedes Perron-Frobenius-Argument, das Mischung (Irreduziblität +
        Aperiodizität) voraussetzt, kann auf $P_{\text{EABC}}$ angewendet
        werden --- \emph{wenn} Mischung gezeigt ist.
  \item Die Existenz eines eindeutigen stationären Maßes~$\mu$ auf
        $\{1,5,7,11\}$ ist äquivalent zu $\rho(P_{\text{EABC}}) < 1$
        (da die Matrix substochastisch ist, ist das stationäre Maß kein
        Wahrscheinlichkeitsmaß, sondern ein sub-invariantes Maß).
\end{enumerate}
\end{proposition}

\begin{tcolorbox}[kernidee]
\textbf{Die eigentliche Rolle der 12er-Zahl:}
Das 12-Kuss-Theorem garantiert, dass die Ikosaeder-Eckstruktur
\emph{maximal, aber nicht überfüllt} ist.
Übertragen: $\ZmZ$ hat genau 12 Elemente --- kein 13.\ Element existiert.
Diese \emph{Vollständigkeit ohne Redundanz} ist der tiefste gemeinsame Kern
von Schütte-Spannung und EABC-Struktur.
Sie liefert:
\begin{itemize}[noitemsep]
  \item Die Endlichkeit des Zustandsraums (Proposition~\ref{prop:endlich}).
  \item Die prinzipielle Anwendbarkeit von Spektralargumenten.
  \item Aber \emph{nicht}: die konkrete Mischungsrate oder den
        Lyapunov-Exponenten.
\end{itemize}
\end{tcolorbox}

% ====================================================
\section{Implikationen für den Lyapunov-Exponenten}
\label{sec:lyapunov}
% ====================================================

\subsection{Rückblick: Die konditionierte Driftrechnung}

Aus \textit{collatz\_lyapunov\_konditioniert.tex} gilt:
\[
  \E[\logF] \;\approx\; -0{,}330 \;<\; 0,
\]
unter der Näherung $\E[\vii \mid l{=}k] \approx 2$ für $k \geq 1$.
Der kritische Schwellenwert ist $\alpha^* \approx 1{,}34$: Für
$\E[\vii \mid l\geq 1] > \alpha^*$ bleibt der Drift negativ.

\subsection{Was die Ikosaeder-Perspektive hinzufügt}

\begin{tcolorbox}[ikosaeder, title={Beitrag 1: Kompaktheitsargument (Analogie)}]
Das Schütte-Theorem ist ein \emph{Kompaktheitssatz}: Es gibt keine
unbeschränkte Folge von Küssvektoren --- die Zahl ist endlich.

Übertragen auf Collatz: Der EABC-Zustandsraum ist endlich (4 Zustände).
Dies bedeutet:
\begin{itemize}[noitemsep]
  \item Die Collatz-Bahn kann nicht ``einfach entweichen'' --- sie muss
        ständig durch den EABC-Filter.
  \item Jede Iteration $\mathcal{U}$ bildet in $\{1, 5, 7, 11\} \pmod{12}$
        ab --- der Zustandsraum hat keine ``freien'' Dimensionen.
\end{itemize}
\textbf{Status:} Diese Kompaktheit ist bewiesen (Endlichkeit von $\ZmZ$),
aber sie liefert keine direkte Schranke für $\E[\logF]$.
\end{tcolorbox}

\begin{tcolorbox}[ikosaeder, title={Beitrag 2: Winkeldefekt als Log-Drift (Analogie)}]
Der Winkeldefekt des Ikosaeders $\delta_{\text{Iko}} = \pi/3$ misst,
wie weit das Ikosaeder von flacher (planarer) Symmetrie abweicht.
Numerisch:
\[
  \delta_{\text{Iko}} / (2\pi) \;=\; \tfrac{1}{6} \;\approx\; 0{,}167.
\]
Der konditionierte Lyapunov-Exponent (in Basis-2-Logarithmen):
\[
  \E[\logF] \;\approx\; -0{,}330 \;\approx\; -2 \cdot 0{,}167.
\]
\textbf{Numerische Beobachtung (spekulativ):}
$\E[\logF] \approx -2 \cdot (\delta_{\text{Iko}} / 2\pi)$.

\medskip
\textbf{Status:} Diese numerische Koinzidenz ist nicht kausal erklärt.
Sie ist eine interessante Beobachtung, aber \emph{kein Beweis}.
\end{tcolorbox}

\begin{tcolorbox}[ikosaeder, title={Beitrag 3: Kissing-Bound als Mischungsschranke (Analogie)}]
Wenn die 12 Kugeln des Kissing-Problems den 12 Restklassen entsprechen,
dann entspricht die Tatsache ``keine 13.\ Kugel passt hinzu'' dem Fakt:

\medskip\noindent
\emph{Es gibt keine 13.\ Restklasse mod~12.}

\medskip
Dies ist trivial wahr, aber strukturell bedeutsam: Die EABC-Dynamik ist
in einem abgeschlossenen, maximalen System --- ohne Overflow.
Wenn man annimmt, dass die 4 aktiven EABC-Klassen die 4-Ecken-Tetaeder-
Unterstruktur des Ikosaeders spiegeln, dann liefert die ``Spannung''
zwischen den 12 Kugeln eine Erklärung, warum die inaktiven 8 Klassen
als Drain fungieren: Sie sind die ``nicht-Tetraeder-Ecken'' im Ikosaeder.

\textbf{Status:} Geometrische Analogie; nicht formalisierbar ohne
explizite Einbettung.
\end{tcolorbox}

\subsection{Numerische Zusammenfassung}

\begin{table}[h!]
\centering
\caption{Gegenüberstellung: Ikosaeder und Collatz-Lyapunov}
\label{tab:vergleich}
\renewcommand{\arraystretch}{1.35}
\begin{tabular}{p{5.5cm} c p{6cm}}
\toprule
\textbf{Ikosaeder / Schütte} & & \textbf{Collatz / EABC} \\
\midrule
12 Ecken & $\longleftrightarrow$ & 12 Restklassen mod 12 \\
4 Tetraeder-Ecken & $\longleftrightarrow$ & 4 EABC-Klassen $\{1,5,7,11\}$ \\
8 restliche Ecken & $\longleftrightarrow$ & 8 inaktive Klassen \\
Kissing-Zahl = 12 & $\longleftrightarrow$ & $|\ZmZ| = 12$ \\
Winkeldefekt $\delta = \pi/3$ & $\longleftrightarrow$ & Drift $\approx -\pi/3 \cdot (1/\pi)$?\\
Schütte-Spannung $\sigma > 0$ & $\longleftrightarrow$ & Heuristische Drift-Lücke \\
Kissing $= 12$, nicht 13 & $\longleftrightarrow$ & Keine 13.\ Restklasse mod 12 \\
\bottomrule
\end{tabular}
\end{table}

% ====================================================
\section{Was lässt sich tatsächlich beweisen? Was bleibt Analogie?}
\label{sec:status}
% ====================================================

\subsection{Vollständige Statusübersicht}

\begin{table}[h!]
\centering
\caption{Status aller Aussagen in diesem Dokument}
\label{tab:status}
\renewcommand{\arraystretch}{1.4}
\begin{tabular}{p{7.5cm} c p{4cm}}
\toprule
\textbf{Aussage} & \textbf{Status} & \textbf{Quelle} \\
\midrule
12-Kuss-Theorem (Schütte/van der Waerden 1953)
  & \textcolor{bewiesengrün}{\textbf{bewiesen}}
  & Klassisches Resultat \\
Ikosaeder hat genau 12 Ecken
  & \textcolor{bewiesengrün}{\textbf{bewiesen}}
  & Euklid / Euler \\
$|\ZmZ| = 12$ Restklassen
  & \textcolor{bewiesengrün}{\textbf{bewiesen}}
  & Trivial \\
EABC-Klassen $= (\Z/12\Z)^\times = \{1,5,7,11\}$
  & \textcolor{bewiesengrün}{\textbf{bewiesen}}
  & Gruppentheorie \\
$P(l = k) = 2^{-(k+1)}$ für $k \geq 1$
  & \textcolor{bewiesengrün}{\textbf{bewiesen}}
  & \textit{coll.\_lyap.\_kond.tex} \\
LTE-Schranke $\vii(3n_l+1) = O(\log l)$
  & \textcolor{bewiesengrün}{\textbf{bewiesen}}
  & LTE (Kummer) \\
$P_{\text{EABC}}$ ist endliche $(4 \times 4)$-Matrix
  & \textcolor{bewiesengrün}{\textbf{bewiesen}}
  & Endlichkeit von $\ZmZ$ \\
$\rho(P_{\text{EABC}}) \leq 1$
  & \textcolor{bewiesengrün}{\textbf{bewiesen}}
  & Perron-Frobenius \\
\midrule
Analogie: Ikosaeder-Ecken $\leftrightarrow$ EABC-Klassen
  & \textcolor{warnorange}{\textbf{Analogie}}
  & Dieses Dokument \\
Numerische Koinzidenz $\E[\logF] \approx -2\delta/(2\pi)$
  & \textcolor{warnorange}{\textbf{spekulativ}}
  & Dieses Dokument \\
Kissing-Bound als Mischungsschranke
  & \textcolor{warnorange}{\textbf{Analogie}}
  & Dieses Dokument \\
$\E[\logF] \approx -0{,}330$
  & \textcolor{warnorange}{\textbf{heuristisch}}
  & \textit{coll.\_lyap.\_kond.tex} \\
\midrule
$\rho(P_{\text{EABC}}) < 1$ (strikt)
  & \textcolor{lückenrot}{\textbf{offen}}
  & Mischung nötig \\
Ergodizität der Collatz-Bahn
  & \textcolor{lückenrot}{\textbf{offen}}
  & Birkhoff fehlt \\
Zyklenfreiheit für alle $n$
  & \textcolor{lückenrot}{\textbf{offen}}
  & Nicht bewiesen \\
Direkter Zusammenhang Schütte $\to$ Lyapunov
  & \textcolor{lückenrot}{\textbf{offen}}
  & Kein formaler Beweis \\
\bottomrule
\end{tabular}
\end{table}

\subsection{Die drei Ebenen im Klartext}

\subsubsection{Ebene 1: Was bewiesen ist}

\begin{tcolorbox}[ergebnis, title={Beweisbare Aussagen}]
\begin{enumerate}[noitemsep]
  \item Das 12-Kuss-Theorem ist ein klassisches Resultat der Geometrie.
  \item Die EABC-Klassen $\{1, 5, 7, 11\} \subset \ZmZ$ sind die
        Einheitengruppe $(\Z/12\Z)^\times$; das ist Gruppentheorie.
  \item Die Übergangsmatrix $P_{\text{EABC}}$ ist eine endliche
        substochastische $(4 \times 4)$-Matrix mit $\rho \leq 1$.
  \item Die LTE-Schranke (Lifting-the-Exponent) liefert
        $\E[\vii \mid l \geq 1] \geq 2$, was den Drift unter dem
        kritischen Wert~$\alpha^* \approx 1{,}34$ hält.
  \item Die Zahl~12 erscheint doppelt: als Kissing-Zahl in $\R^3$
        und als Modulus des EABC-Rahmens. Diese numerische Koinzidenz
        ist mathematisch real, aber kausal nicht erklärt.
\end{enumerate}
\end{tcolorbox}

\subsubsection{Ebene 2: Geometrische Analogie}

\begin{tcolorbox}[analogie]
\textbf{Was die Analogie strukturell zeigt (aber nicht beweist):}
\begin{itemize}[noitemsep]
  \item Die EABC-Dynamik lebt in einem \emph{maximal kompakten},
        endlichen Zustandsraum --- analog zur Maximalität der 12 Kissing-Kugeln.
  \item Die ``Spannung'' im Ikosaeder (kein starres Gleichgewicht,
        Gregory-Newton-Lücke) spiegelt die ``heuristische Lücke'' in der
        Collatz-Drift-Rechnung: Alles deutet auf negativen Drift, aber
        keiner hat Ergodizität bewiesen.
  \item Die 4-Teilung EABC innerhalb der 12-Teilung mod~12 entspricht
        geometrisch dem Einbetten eines Tetraeders in das Ikosaeder ---
        der ``aktive Kern'' innerhalb der maximalen Konfiguration.
\end{itemize}
\end{tcolorbox}

\subsubsection{Ebene 3: Was offen bleibt}

\begin{tcolorbox}[lücke, title={Fundamental offene Fragen}]
\begin{enumerate}[noitemsep]
  \item \textbf{Mischung:} Ist $P_{\text{EABC}}$ irreduzibel und aperiodisch?
        Dann wäre $\rho(P_{\text{EABC}}) < 1$ automatisch, und ein
        stationäres Maß existierte --- aber Collatz-Dynamik mit beliebigem
        Startpunkt bleibt komplizierter als eine endliche Markovkette.

  \item \textbf{Ergodizität:} Ohne Birkhoff-Ergodensatz für die
        Collatz-Abbildung~$\mathcal{U}$ folgt aus $\E[\logF] < 0$ nur,
        dass der Drift für \emph{zufällig gewählte} Bahnen negativ ist,
        nicht für alle $n \in \N$.

  \item \textbf{Quantitative Brücke:} Es gibt keine formale Übertragung
        der Schütte-Ungleichung (Lemma~\ref{lem:schuette}) auf eine
        Spektralschranke für $P_{\text{EABC}}$.
        Die Matrixeinträge sind Übergangswahrscheinlichkeiten, keine
        Winkelkosinusse von Einheitsvektoren.

  \item \textbf{Winkeldefekt-Drift:} Die numerische Koinzidenz
        $\E[\logF] \approx -2\delta_{\text{Iko}}/(2\pi)$ ist ungeklärt.
        Sie könnte Zufall sein.
\end{enumerate}
\end{tcolorbox}

\subsection{Präzisiertes Hauptresultat dieses Dokuments}

\begin{tcolorbox}[hauptsatz, title={Stärkster beweisbarer Satz (dieses Dokuments)}]
Die Zahl~$12$ erscheint in zwei unabhängigen mathematischen Kontexten:
\begin{enumerate}[noitemsep]
  \item Als Kissing-Zahl in $\R^3$ (Schütte/van der Waerden, 1953):
        maximal viele, geometrisch gespannte Kugeln ohne Überlappung.
  \item Als Modulus des EABC-Rahmens: $(\Z/12\Z)^\times = \{1, 5, 7, 11\}$,
        mit $|(\Z/12\Z)^\times| = 4$ aktiven und $8$ inaktiven Klassen.
\end{enumerate}

\medskip
Beide Strukturen sind maximal kompakt und endlich.
Aus der Endlichkeit folgt:
\begin{itemize}[noitemsep]
  \item $P_{\text{EABC}}$ ist eine endliche Matrix (Proposition~\ref{prop:endlich}).
  \item Spektralargumente sind \emph{prinzipiell anwendbar}, sobald Mischung
        gezeigt ist.
  \item Die negative konditionierte Drift
        $\E[\logF] \approx -0{,}330$ (aus
        \textit{collatz\_lyapunov\_konditioniert.tex}) ist innerhalb
        des endlichen Zustandsraums konsistent.
\end{itemize}

\medskip
\textbf{Was dies \emph{nicht} beweist:}
Die Collatz-Vermutung.
Die Ikosaeder-Analogie verstärkt die geometrische Kohärenz des Rahmens,
ersetzt aber nicht die fehlenden Schritte: Ergodizität, Zyklenfreiheit,
und einen quantitativen Beweis der Drift-Negativität für alle Bahnen.
\end{tcolorbox}

% ====================================================
\appendix
\section{Ikosaeder-Daten und Adjazenzstruktur}
\label{app:iko}
% ====================================================

Das reguläre Ikosaeder $\mathcal{I}$ hat Koordinaten der 12 Ecken
(in Standardlage):
\[
  V(\mathcal{I}) \;=\; \bigl\{
    (\pm 1, 0, \pm\varphi),\;
    (0, \pm\varphi, \pm 1),\;
    (\pm\varphi, \pm 1, 0)
  \bigr\},
\]
mit dem goldenen Schnitt $\varphi = (1+\sqrt{5})/2$.
Die Gram-Identität der Adjazenzmatrix (vgl.\ \textit{IkoDo\_transition.tex}):
\[
  B_I^\top B_I \;=\; 5\,I_{12} + A_I,
\]
denn jede Ikosaederecke liegt auf genau 5 Kanten.

Der Winkel zwischen benachbarten Ecken:
\[
  \cos\theta_{\text{benachb}} \;=\;
  \frac{\langle v_i, v_j\rangle}{\|v_i\|\|v_j\|}
  \;=\; \frac{1}{\sqrt{5}} \;\approx\; 0{,}447,
  \qquad \theta \;\approx\; 63{,}4°.
\]
Dieser Winkel liegt knapp über dem Kissing-Mindestwinkel $60°$, was die
``Spannung'' illustriert: Die Ikosaeder-Ecken sind fast, aber nicht ganz
auf der Grenze der Kissing-Bedingung.

\section{Verbindung zu den Vorgängerdokumenten}
\label{app:verbindung}

Dieses Dokument steht in folgender Beziehung zu den anderen Collatz-Texten
des Projekts:

\begin{center}
\renewcommand{\arraystretch}{1.3}
\begin{tabular}{l p{9cm}}
\toprule
\textbf{Datei} & \textbf{Inhalt und Verbindung} \\
\midrule
\textit{collatz\_lyapunov.tex}
  & Naive Lyapunov-Analyse, $\E[\logF] \approx -0{,}83$
    (geometrische $\nu$-Verteilung, unkonditioniert) \\
\textit{collatz\_lyapunov\_konditioniert.tex}
  & Konditionierte Analyse, $\E[\logF] \approx -0{,}33$;
    LTE-Präzisierung; Sensitivitätsanalyse;
    kritischer Wert $\alpha^* \approx 1{,}34$ \\
\textit{collatz\_eabc\_spektrum.tex}
  & EABC-Klassenstruktur; Spannungs-/Drain-Regime \\
\textit{collatz\_ikosaeder\_spannung.tex}
  & Dieses Dokument: geometrische Analogie Ikosaeder $\leftrightarrow$ EABC \\
\bottomrule
\end{tabular}
\end{center}

\end{document}
